\documentclass[12pt]{article}
\usepackage{amsmath,amssymb,amsthm}
\usepackage[margin=1in]{geometry}

\newtheorem{theorem}{Theorem}
\newtheorem{lemma}[theorem]{Lemma}

\title{Problem 193: Squarefree Numbers}
\author{Project Euler}
\date{}

\begin{document}
\maketitle

\section{Problem Statement}

Count the squarefree numbers below $2^{50}$.

\section{Mathematical Foundation}

\begin{theorem}[Squarefree Indicator Identity]
For all positive integers $n$,
\[
|\mu(n)|^2 = \sum_{d^2 \mid n} \mu(d).
\]
\end{theorem}

\begin{proof}
Both sides are multiplicative, so it suffices to check at prime powers $n = p^k$:
\begin{itemize}
    \item $k = 1$: $|\mu(p)|^2 = 1$ and $\sum_{d^2 \mid p} \mu(d) = \mu(1) = 1$. \checkmark
    \item $k \geq 2$: $|\mu(p^k)|^2 = 0$ and $\sum_{d^2 \mid p^k} \mu(d) = \mu(1) + \mu(p) = 1 - 1 = 0$. \checkmark
\end{itemize}
\qed
\end{proof}

\begin{theorem}[Squarefree Counting Formula]
The number of squarefree integers not exceeding $N$ is
\[
Q(N) = \sum_{k=1}^{\lfloor \sqrt{N} \rfloor} \mu(k) \left\lfloor \frac{N}{k^2} \right\rfloor.
\]
\end{theorem}

\begin{proof}
\[
Q(N) = \sum_{n=1}^{N} |\mu(n)|^2 = \sum_{n=1}^{N} \sum_{d^2 \mid n} \mu(d) = \sum_{d=1}^{\lfloor\sqrt{N}\rfloor} \mu(d) \left\lfloor \frac{N}{d^2} \right\rfloor.
\]
The exchange of summation is valid since $d^2 \mid n$ and $n \leq N$ imply $d \leq \sqrt{N}$. \qed
\end{proof}

\begin{lemma}[M\"obius Sieve]
The values $\mu(k)$ for $1 \leq k \leq M$ can be computed in $O(M \log\log M)$ time and $O(M)$ space via a modified Eratosthenes sieve.
\end{lemma}

\begin{proof}
Initialize $\mu(k) = 1$. For each prime $p$: multiply $\mu$ by $-1$ at multiples of $p$; set $\mu = 0$ at multiples of $p^2$. Total work is $\sum_{p \leq M} M/p = O(M \log\log M)$ by Mertens' theorem. \qed
\end{proof}

\section{Algorithm}

\begin{verbatim}
function count_squarefree(N):
    M = floor(sqrt(N))
    Sieve mu(1), ..., mu(M)
    return sum of mu(k) * floor(N / k^2) for k = 1 to M
\end{verbatim}

\section{Complexity Analysis}

\begin{itemize}
    \item \textbf{Time:} $O(\sqrt{N} \log\log\sqrt{N})$ for sieving, $O(\sqrt{N})$ for summation.
    \item \textbf{Space:} $O(\sqrt{N}) \approx 3.4 \times 10^7$ integers for $N = 2^{50}$.
\end{itemize}

\section{Answer}

\[
\boxed{684465067343069}
\]

\end{document}
